% Part: lambda-calculus
% Chapter: church-rosser
% Section: definitions-and-properties

\documentclass[../../../include/open-logic-section]{subfiles}

\begin{document}

\olfileid{lam}{cr}{dap}

\olsection{વ્યાખ્યા અને ગુણધર્મો}

આ અધ્યાયમાં ચર્ચ--રોસર ગુણધર્મનો ખ્યાલ અને તેના
કેટલાક સામાન્ય ગુણધર્મો રજૂ કરીએ છીએ.

\begin{defn}[ચર્ચ--રોસર ગુણધર્મ, CR]
  પદો પરનો સંબંધ~$\xredone$ \emph{ચર્ચ--રોસર ગુણધર્મ}
  ધરાવે છે ત્યારે અને માત્ર ત્યારે જ જ્યારે
  $M \xredone P$ અને $M \xredone Q$ હોય ત્યારે
  એવું~$N$ મળે કે $P \xredone N$ અને
  $Q \xredone N$.
\end{defn}

લૅમ્બડા કલનને સંગણનાના નમૂના તરીકે જોઈ શકાય: તેમાં
સામાન્ય સ્વરૂપનાં પદો ``મૂલ્યો'' છે અને પદો પરનો
લઘુકરણ-સંબંધ ``ગણતરીના નિયમો'' આપે છે. ચર્ચ--રોસર
ગુણધર્મ કહે છે કે ગણતરીમાં આગળ વધવાની એક કરતાં વધુ
રીતો હોય તોપણ અભિવ્યક્તિનું અંતિમ મૂલ્ય એક જ રહે છે.

પ્રાથમિક બીજગણિતનું ઉદાહરણ લઈએ. $4 \times (1+2) + 3$
ગણવાની એકથી વધુ રીતો છે. પહેલા $1+2$ને~$3$ કરીએ તો
$4 \times 3+3$ મળે; અથવા વિતરણ-નિયમથી પહેલા
$4 \times (1+2)$ને વિસ્તારી $4 \times 1+4
\times 2+3$ મળે. જોકે બંનેને આગળ ગણતાં $12+3$ જ
મળે છે.

$\xredone$ને $\beta$-લઘુકરણ માનીએ તો ચર્ચ--રોસર
ગુણધર્મનું પરિણામ છે કે પદને સામાન્ય સ્વરૂપ મળે તો તે
અનન્ય છે. માનો કે~$M$નું $P$ અને $Q$ બંનેમાં લઘુકરણ
થાય છે અને બંને સામાન્ય સ્વરૂપો છે. ચર્ચ--રોસર
ગુણધર્મથી એવું~$N$ મળે કે $P$ અને $Q$ બંનેનું
તેમાં લઘુકરણ થાય. $P$ અને $Q$ સામાન્ય સ્વરૂપમાં
હોવાથી $P$ અને $Q$નું $N$માં લઘુકરણ ફક્ત તુચ્છ,
એટલે શૂન્ય-પગલાનું જ
હોઈ શકે; તેથી $P$, $Q$ અને $N$ એક જ પદ છે. આથી
કોઈ પદના \emph{તે} સામાન્ય સ્વરૂપની વાત કરવી યોગ્ય છે.

લૅમ્બડા કલનને સંગણનાના નમૂના તરીકે જોતાં પદનું સામાન્ય
સ્વરૂપ તે પદથી શરૂ થતી ગણતરીનું ``અંતિમ પરિણામ''
ગણી શકાય. ઉપરના ઉપસિદ્ધાંત મુજબ આવું અંતિમ પરિણામ
હોય તો માત્ર એક જ હોય છે, જેમ કે $4 \times (1+2)+3$
ગણતાં અંતે માત્ર~$15$ મળે છે.

\begin{thm} \ollabel{thm:str}
  સંબંધ~$\xredone$ ચર્ચ--રોસર ગુણધર્મ ધરાવતો હોય
  અને $\xred$ એ $\xredone$ને સમાવતો સૌથી નાનો
  પરંપરિત સંબંધ હોય, તો $\xred$ પણ ચર્ચ--રોસર
  ગુણધર્મ ધરાવે છે.
\end{thm}

\begin{proof}
  માનો કે
  \begin{align*}
    M & \xredone P_1 \xredone \dots \xredone P_m \text{ અને}\\
    M & \xredone Q_1 \xredone \dots \xredone Q_n.
  \end{align*}
  ઊંચાઈ~$m + 1$ અને પહોળાઈ~$n + 1$ ધરાવતું
  પદોનું જાળક~$N$ રચીને પ્રમેય સાબિત કરીશું.
  $N_{i,j}$ એ $i$મી પંક્તિ અને $j$મા સ્તંભનું પદ છે.
  
  જાળક~$N$ એવું રચીએ કે $N_{i,j} \xredone N_{i+1,j}$
  અને $N_{i,j} \xredone N_{i,j+1}$. તેની વ્યાખ્યા:
  \begin{align*}
    N_{0,0} &= M \\
    N_{i,0} &= P_i && \text{જો } 1 \le i \le m \\
    N_{0,j} &= Q_j && \text{જો } 1 \le j \le n \\
  \intertext{અને બાકી કિસ્સામાં:}
    N_{i,j} &= R
  \end{align*}
  અહીં~$R$ એવું પદ છે કે $N_{i-1,j} \xredone R$ અને
  $N_{i,j-1} \xredone R$. $\xredone$ના ચર્ચ--રોસર
  ગુણધર્મથી આવું પદ હંમેશા મળે છે.

  હવે $N_{m,0} \xredone \dots \xredone N_{m,n}$
  અને $N_{0,n} \xredone \dots \xredone N_{m,n}$.
  અહીં $N_{m,0}$ એ $P_m$ અને $N_{0,n}$ એ $Q_n$ છે.
  $\xred$ની વ્યાખ્યાથી દાવો મળે છે.
  \sourcecorrection{OLLAM-032}{સ્થિર અંગ્રેજી મૂળમાં
  સાબિતીના અંતે $N_{m,0}$ને $P$ અને $N_{0,n}$ને
  $Q$ કહે છે, પરંતુ આ સાબિતીમાં ફક્ત~$P_i$ અને
  $Q_j$ વ્યાખ્યાયિત છે. અહીં ચોક્કસ અંતિમ પદો
  $P_m$ અને $Q_n$ લખ્યાં છે.}
\end{proof}

\end{document}
